\begin{lemma}[Base invariance of the bqo restriction property]
Assume that $r$ is a quasi-order and $Q$ is bqo.  For every infinite base
$X$ and every locally constant $h:\mathsf{BaseIncSeq}(X)\to Q$, there is
an enumeration $z\in\mathsf{BaseIncSeq}(X)$ such that the global-coordinate
restriction $x\mapsto h(z\circ x)$ is perfect.
\end{lemma}
\begin{proof}
Applying \noderef{InfiniteSetEnumeration} to $X$ and the given infinitude
$hX$, choose $e\in\mathsf{IncSeq}$ with
\[
\operatorname{range}(e)=X.
\]
For each $x\in\mathsf{IncSeq}$, \noderef{IncSeqComp} gives an increasing
embedding $\mathsf{IncSeqComp}(e,x)$ satisfying
\[
\mathsf{IncSeqComp}(e,x)(i)=e(x(i))
\quad\text{for every }i.
\]
Its range is contained in $X$: if
$n=\mathsf{IncSeqComp}(e,x)(i)=e(x(i))$, then
$n\in\operatorname{range}(e)=X$.  Hence, by the range condition in
\noderef{BaseIncSeq},
\[
\widehat x=\bigl\langle\mathsf{IncSeqComp}(e,x),
  \operatorname{range}(\mathsf{IncSeqComp}(e,x))\subseteq X\bigr\rangle
\in\mathsf{BaseIncSeq}(X).
\]
The domain assertion in \noderef{BaseMultiSequence} says that $h$ accepts
exactly such base increasing enumerations.  Define the pullback
$\widetilde h:\mathsf{IncSeq}\to Q$ by
\[
\widetilde h(x)=h(\widehat x).
\]
Thus $\widetilde h$ is a multi-sequence in the sense of
\noderef{MultiSequence}.

We next verify local constancy of $\widetilde h$ in full detail.  Fix
$x\in\mathsf{IncSeq}$ and apply \noderef{BaseLocallyConstant} to the base
enumeration $\widehat x$.  There is an $n\in\mathbb N$ such that for every
$b\in\mathsf{BaseIncSeq}(X)$,
\[
\mathsf{PrefixAgree}(n,\widehat x.1,b.1)
\quad\Longrightarrow\quad h(b)=h(\widehat x).
\tag{1}
\]
Here \noderef{PrefixAgree} means equality in every coordinate $i<n$.
Let $y\in\mathsf{IncSeq}$ satisfy
$\mathsf{PrefixAgree}(n,x,y)$, and form the base enumeration
\[
\widehat y=\bigl\langle\mathsf{IncSeqComp}(e,y),
  \operatorname{range}(\mathsf{IncSeqComp}(e,y))\subseteq X\bigr\rangle,
\]
which is in $\mathsf{BaseIncSeq}(X)$ by the same
\noderef{BaseIncSeq} and range argument as above.  If $i<n$, then
$x(i)=y(i)$ by \noderef{PrefixAgree}; applying $e$ to this equality and
using \noderef{IncSeqComp} gives
\[
\mathsf{IncSeqComp}(e,x)(i)=e(x(i))=e(y(i))
  =\mathsf{IncSeqComp}(e,y)(i).
\]
Therefore
$\mathsf{PrefixAgree}(n,\widehat x.1,\widehat y.1)$, and (1) gives
\[
\widetilde h(y)=h(\widehat y)=h(\widehat x)=\widetilde h(x).
\]
Since $x$ was arbitrary, this is precisely the definition of
\noderef{LocallyConstant}; hence $\widetilde h$ is locally constant.

Apply \noderef{ShiftDichotomy} to $\widetilde h$.  We obtain a
$u\in\mathsf{IncSeq}$ such that either
\[
\mathsf{PerfectMultiSequence}\,r
  (\mathsf{MultiSequenceRestrict}\,\widetilde h\,u)
\tag{2}
\]
or
\[
\mathsf{BadMultiSequence}\,r
  (\mathsf{MultiSequenceRestrict}\,\widetilde h\,u).
\tag{3}
\]
The predicate displayed in (3) is the badness predicate of
\noderef{BadMultiSequence}.  Suppose (3) holds.  By
\noderef{RestrictionLocallyConstant}, the local
constancy of $\widetilde h$ implies
\[
\mathsf{LocallyConstant}\bigl(\mathsf{MultiSequenceRestrict}\,\widetilde h\,u\bigr).
\]
But \noderef{Bqo} says that the hypothesis $hb$ rules out every bad locally
constant multi-sequence in $Q$: applying it to the restricted multi-sequence
gives the negation of (3), a contradiction.  Consequently (2) holds.

Set
\[
z=\mathsf{IncSeqComp}(e,u).
\]
By \noderef{IncSeqComp}, $z\in\mathsf{IncSeq}$.  If
$n=z(i)=e(u(i))$, then $n\in\operatorname{range}(e)=X$; hence
$\operatorname{range}(z)\subseteq X$, and the definition of
\noderef{BaseIncSeq} gives
\[
z_X=\langle z,\operatorname{range}(z)\subseteq X\rangle
\in\mathsf{BaseIncSeq}(X).
\]

It remains to identify the two restrictions pointwise.  Fix
$x\in\mathsf{IncSeq}$.  By \noderef{BaseRestrict}, the value
$(\mathsf{BaseRestrict}(h,z_X))(x)$ is $h$ applied to the base enumeration
whose underlying increasing embedding is $\mathsf{IncSeqComp}(z,x)$; its
range-in-$X$ proof is exactly the proof supplied in the definition of
\noderef{BaseRestrict}.  For every coordinate $i$,
\noderef{IncSeqComp} gives
\[
\mathsf{IncSeqComp}(z,x)(i)=z(x(i))=e(u(x(i)))
 =\mathsf{IncSeqComp}(e,\mathsf{IncSeqComp}(u,x))(i).
\tag{4}
\]
Thus these two increasing embeddings are equal pointwise, and the
corresponding base-enumeration subtype values are equal.  Using
\noderef{MultiSequenceRestrict} and the definition of $\widetilde h$, we
therefore obtain
\[
\begin{aligned}
(\mathsf{BaseRestrict}(h,z_X))(x)
 &=h\bigl(\bigl\langle\mathsf{IncSeqComp}(e,\mathsf{IncSeqComp}(u,x)),
     \operatorname{range}(\mathsf{IncSeqComp}(e,\mathsf{IncSeqComp}(u,x)))
       \subseteq X\bigr\rangle\bigr)\\
 &=\widetilde h(\mathsf{IncSeqComp}(u,x))\\
 &=(\mathsf{MultiSequenceRestrict}\,\widetilde h\,u)(x).
\end{aligned}
\tag{5}
\]

Finally, \noderef{PerfectMultiSequence} states that (2) means, for every
$x\in\mathsf{IncSeq}$,
\[
r\bigl((\mathsf{MultiSequenceRestrict}\,\widetilde h\,u)(x),
  (\mathsf{MultiSequenceRestrict}\,\widetilde h\,u)(\mathsf{ShiftMap}(x))\bigr).
\]
Apply (5) once to $x$ and once to $\mathsf{ShiftMap}(x)$.  The resulting
statement is
\[
r\bigl((\mathsf{BaseRestrict}(h,z_X))(x),
  (\mathsf{BaseRestrict}(h,z_X))(\mathsf{ShiftMap}(x))\bigr)
\]
for every $x$, which is exactly the defining universal condition in
\noderef{PerfectMultiSequence}.  Hence
$\mathsf{PerfectMultiSequence}\,r(\mathsf{BaseRestrict}(h,z_X))$, and
$z_X$ is the required element of $\mathsf{BaseIncSeq}(X)$; renaming this
element as $z$ gives the enumeration in the lemma statement.
\end{proof}
